\begin{proof}[Proof of \cref{obj:thm:full-honest-frontier-answer}]
\begin{enumerate}
\item Choose the constants \(c,C\) supplied by
\cref{obj:thm:full-honest-frontier-solution}, with
\[
0<c,\qquad 0<C.
\]
Its universal construction guarantee applies simultaneously to every admissible engine and naming policy, every \((\alpha,\beta)\in\mathcal D\), every integer \(n\ge1\), and every \(r\in[0,1/2]\), with these same constants. Its compact-uniform guarantee, elbow identities, and converse for every globally honest procedure, including independently randomized procedures, give the corresponding assertions of the present theorem. In particular, its choices
\[
c_\star=c,\qquad C_\star=C,\qquad n_0=1
\]
give, for every nonempty compact \(K\subseteq\mathcal D\),
\[
0<\inf_{(\alpha,\beta)\in K}c_\star,\qquad
\sup_{(\alpha,\beta)\in K}C_\star<\infty,\qquad
\sup_{(\alpha,\beta)\in K}n_0=1.
\]
% lean: CausalSmith.Stat.LogoddsLowsmoothFrontier.full_honest_frontier_solution

\item Fix \((\alpha,\beta)\in\mathcal D\) and an integer \(n\ge1\). The concrete sequence in
\cref{obj:thm:full-honest-frontier-solution} is exactly
\[
I_n^\star(\alpha,\beta)
=I^{\mathrm{up}}_{n,\mathsf E_0,\mathsf N_0}(\alpha,\beta),
\]
as specified in \cref{obj:def:frontier-handle}. That guarantee yields
\[
I_n^\star(\alpha,\beta)\in\mathcal A_n(\alpha,\beta).
\]
For every sample and seed, it also gives rational endpoints \(l,u\) such that
\[
I_n^\star(\alpha,\beta)(\mathcal O,U)=[l,u],
\qquad
-\frac12\le l\le u\le\frac12.
\]
% lean: CausalSmith.Stat.LogoddsLowsmoothFrontier.full_honest_frontier_solution

\item Fix \(r\in[0,1/2]\). Expanding the diagnostic envelope of
\cref{obj:def:diagnostic-envelope} and its identification with the rate profile in
\cref{obj:def:frontier-handle} gives
\[
\rho_n(\alpha,\beta;r)=d_n(\alpha,\beta;r)
=
n^{-\min\{1/2,\,2(\alpha+\beta)/(2\alpha+2\beta+1)\}}
+r\,n^{-4\beta/(4\beta+1)}.
\]
Thus the concrete comparison supplied by
\cref{obj:thm:full-honest-frontier-solution}, expressed with this explicit profile, is
\[
c\,\rho_n(\alpha,\beta;r)
\le J_n(\alpha,\beta;r)
\le
\sup_{P\in\mathcal M_r(\alpha,\beta)}
E_{P^{\otimes n}\otimes\lambda}
\bigl|I_n^\star(\alpha,\beta)(\mathcal O,U)\bigr|
\le C\,\rho_n(\alpha,\beta;r).
\]
Together with the assertions retained in the first step, this establishes all the stated guarantees with the same constants.
% lean: CausalSmith.Stat.LogoddsLowsmoothFrontier.full_honest_frontier_answer
\end{enumerate}
\end{proof}

\begin{proof}[Proof of \cref{obj:thm:finite-honest-upper}]
We first establish the bounds for \(n\ge2\). Fix an admissible arithmetic engine \(\mathsf E\), an admissible naming policy \(\mathsf N\), and public exponents satisfying \(0<\beta<1/4\) and \(\beta<\alpha\le1\).

\begin{enumerate}
\item Define the retained ranks, smoothness sum, and rate exponents by
\[
\begin{gathered}
k_C=k_C(n,\alpha,\beta;\mathsf E,\mathsf N),
\qquad
k_S=k_S(n,\beta;\mathsf E,\mathsf N),
\qquad s=\alpha+\beta,\\
\mathsf a(\alpha,\beta)
=\min\left\{\frac12,\frac{2s}{2s+1}\right\},
\qquad
\mathsf b(\beta)=\frac{4\beta}{4\beta+1}.
\end{gathered}
\]
These are the ranks and exponents of
\cref{obj:def:resolutions,obj:def:diagnostic-envelope}.
For \(k\ge1\), put
\[
V_n(k)=\frac{10}{n}+\frac{4k}{n(n-1)},
\qquad
b_C=8k_C^{-s}+\sqrt{20V_n(k_C)},
\qquad
b_S=25k_S^{-2\beta}+\sqrt{20V_n(k_S)}.
\]
Here \(b_C,b_S\) are the coordinate error radii, as in
\cref{obj:def:upper-handle}.

By \cref{obj:lem:fixed-budget-realization}, with target resolutions
\[
R_C=n^{2/(2s+1)},
\qquad R_S=n^{2/(4\beta+1)},
\]
the ranks satisfy
\[
1\le R_C\le k_C\le3R_C,
\qquad
1\le R_S\le k_S\le3R_S.
\]
In particular, both ranks are positive and are fixed before sampling.

For either pair
\[
(\gamma_{\mathrm{bal}},k_{\mathrm{bal}})
\in\{(s,k_C),(2\beta,k_S)\},
\]
the lower rank bound and \(\gamma_{\mathrm{bal}}>0\) give
\[
k_{\mathrm{bal}}^{-\gamma_{\mathrm{bal}}}
\le
n^{-2\gamma_{\mathrm{bal}}/(2\gamma_{\mathrm{bal}}+1)}.
\]
Also \(n(n-1)\ge n^2/2>0\), so the upper rank bound yields
\[
20V_n(k_{\mathrm{bal}})
\le
\frac{200}{n}
+
480\,\frac{n^{2/(2\gamma_{\mathrm{bal}}+1)}}{n^2}
=
200(n^{-1/2})^2
+
480\left(
n^{-2\gamma_{\mathrm{bal}}/(2\gamma_{\mathrm{bal}}+1)}
\right)^2.
\]
Both powers are nonnegative. The last expression is bounded by
\[
\left[
40\left(
n^{-1/2}
+
n^{-2\gamma_{\mathrm{bal}}/(2\gamma_{\mathrm{bal}}+1)}
\right)
\right]^2,
\]
because expansion of this square gives coefficients \(1600\) on both squared terms and a nonnegative cross term. Taking square roots therefore gives
\[
\sqrt{20V_n(k_{\mathrm{bal}})}
\le
40\left(
n^{-1/2}
+
n^{-2\gamma_{\mathrm{bal}}/(2\gamma_{\mathrm{bal}}+1)}
\right).
\]
Each power on the right is at most the power with exponent
\(-\min\{1/2,2\gamma_{\mathrm{bal}}/(2\gamma_{\mathrm{bal}}+1)\}\).
Since \(0<\beta<1/4\),
\[
\frac{4\beta}{4\beta+1}<\frac12.
\]
Consequently,
\[
\begin{aligned}
b_C
&\le88\,n^{-\mathsf a(\alpha,\beta)}
\le128\,n^{-\mathsf a(\alpha,\beta)},\\
b_S
&\le105\,n^{-\mathsf b(\beta)}
\le128\,n^{-\mathsf b(\beta)}.
\end{aligned}
\]
The exponents are positive, and \(\mathsf a(\alpha,\beta)\le1/2\le1\). Thus
\[
n^{-\mathsf a(\alpha,\beta)}\le1,
\qquad
n^{-\mathsf b(\beta)}\le1,
\qquad
\frac1n\le n^{-\mathsf a(\alpha,\beta)}.
\]
% lean: CausalSmith.Stat.LogoddsLowsmoothFrontier.coordinate_radii_power_bounds
% lean: CausalSmith.Stat.LogoddsLowsmoothFrontier.diagnostic_power_bounds

\item With the denominator floor and sensitivity coefficients defined by
\[
s_0=\frac1{512},
\qquad
\mathsf L_{\mathrm{sens}}=\frac{2\exp(1/2)}{s_0},
\qquad
\mathsf H_{\mathrm{sens}}=\frac{2\exp(1/2)}{s_0^2},
\]
choose
\[
K_0
=
128\mathsf L_{\mathrm{sens}}
+32768\mathsf H_{\mathrm{sens}}
+128\mathsf H_{\mathrm{sens}}\exp(1/2)+2.
\]
All its summands are nonnegative, so \(K_0\ge1>0\).
For every \(r\in[0,1/2]\), the preceding bounds imply
\[
\begin{aligned}
&\mathsf L_{\mathrm{sens}}b_C
+\mathsf H_{\mathrm{sens}}b_S
   \{\exp(1/2)r+2b_C\}
+\frac2n\\
&\quad\le
128\mathsf L_{\mathrm{sens}}n^{-\mathsf a(\alpha,\beta)}
+128\mathsf H_{\mathrm{sens}}n^{-\mathsf b(\beta)}
 \{\exp(1/2)r+256n^{-\mathsf a(\alpha,\beta)}\}
+2n^{-\mathsf a(\alpha,\beta)}\\
&\quad\le
(128\mathsf L_{\mathrm{sens}}
 +32768\mathsf H_{\mathrm{sens}}+2)
 n^{-\mathsf a(\alpha,\beta)}
+
128\mathsf H_{\mathrm{sens}}\exp(1/2)\,
r n^{-\mathsf b(\beta)}\\
&\quad\le
K_0\left\{
n^{-\mathsf a(\alpha,\beta)}
+r n^{-\mathsf b(\beta)}
\right\}
=
K_0d_n(\alpha,\beta;r).
\end{aligned}
\]
The second inequality uses
\(n^{-\mathsf b(\beta)}n^{-\mathsf a(\alpha,\beta)}
\le n^{-\mathsf a(\alpha,\beta)}\).
The last inequality follows by adding the nonnegative terms that complete the common coefficient \(K_0\).
This choice depends solely on the fixed denominator floor.
% lean: CausalSmith.Stat.LogoddsLowsmoothFrontier.radius_sensitivity_rate_assembly
% lean: CausalSmith.Stat.LogoddsLowsmoothFrontier.upper_radii_uniform_rate

\item Fix \(P\in\mathcal M(\alpha,\beta)\). Write
\[
\mathbb P_{P,n}=P^{\otimes n}\otimes\lambda,
\qquad
\mathbb E_{P,n}[\cdot]=\int(\cdot)\,d\mathbb P_{P,n},
\]
where \(\lambda\) is the uniform law of the independent seed. Variances below are under this same probability law.

We first control the unconditional absolute mean of the numerator estimate. Define, for \(x\in[0,1]\),
\[
\begin{gathered}
e(P,x)=\ell\{g(P,x)\},\\
\mu_0(P,x)=\ell\{\nu(P,x)\},
\qquad
\mu_1(P,x)=\ell\{\nu(P,x)+\theta(P)\},\\
m(P,x)=E_P[Y\mid X=x]
=(1-e(P,x))\mu_0(P,x)+e(P,x)\mu_1(P,x).
\end{gathered}
\]
The homogeneous relation and native-logit bounds are those of
\cref{obj:def:model}. The logistic derivative satisfies
\[
0<\ell'(t_{\ell})
=\ell(t_{\ell})\{1-\ell(t_{\ell})\}
\le\frac14
\qquad(t_{\ell}\in\mathbb R).
\]
The derivative bound on the interval between any two real arguments gives
\[
|\ell(t_{\ell})-\ell(t_{\ell}')|
\le\frac14|t_{\ell}-t_{\ell}'|
\qquad(t_{\ell},t_{\ell}'\in\mathbb R).
\]
Hence, for \(x,z\in[0,1]\),
\[
\begin{aligned}
|e(P,x)-e(P,z)|
&\le\frac12|x-z|^\alpha
\le\frac12|x-z|^\beta,\\
|\mu_0(P,x)-\mu_0(P,z)|
&\le\frac12|x-z|^\beta,\\
|\mu_1(P,x)-\mu_1(P,z)|
&\le\frac12|x-z|^\beta,\\
|\mu_1(P,z)-\mu_0(P,z)|
&\le\frac14|\theta(P)|\le\frac18.
\end{aligned}
\]
The comparison of the two covariate powers uses
\(|x-z|\le1\) and \(\beta<\alpha\); it also holds at \(x=z\), since both exponents are positive.

The marginal mean has the decomposition
\[
\begin{aligned}
m(P,x)-m(P,z)
={}&(1-e(P,x))\{\mu_0(P,x)-\mu_0(P,z)\}\\
&+e(P,x)\{\mu_1(P,x)-\mu_1(P,z)\}\\
&+\{e(P,x)-e(P,z)\}
  \{\mu_1(P,z)-\mu_0(P,z)\}.
\end{aligned}
\]
Since \(0<e(P,x)<1\), the preceding moduli imply
\[
\begin{aligned}
|m(P,x)-m(P,z)|
&\le
\{1-e(P,x)\}\frac12|x-z|^\beta
+e(P,x)\frac12|x-z|^\beta
+\frac12|x-z|^\beta\frac18\\
&=\frac9{16}|x-z|^\beta.
\end{aligned}
\]
Thus
\[
[e(P,\cdot)]_\alpha\le\frac12,
\qquad
[m(P,\cdot)]_\beta\le\frac9{16}.
\]
These functions are continuous: the propensity and marginal mean are sums of continuous conditional cells.

Define the empirical treated-outcome mean by
\[
\overline{AY}_n=\frac1n\sum_{i=1}^n A_iY_i.
\]
The coordinate definition in \cref{obj:def:coordinate-estimates} gives
\[
\widehat C_n
=\overline{AY}_n-\mathbb U_{n,k_C}(A,Y),
\qquad
\mathbb E_{P,n}\overline{AY}_n=E_P[AY].
\]
Using the projection expectation and residual identities from
\cref{obj:lem:projection-control}, we obtain
\[
\mathbb E_{P,n}\widehat C_n-C(P)
=
\left\langle
e(P,\cdot)-\Pi_{k_C}e(P,\cdot),
m(P,\cdot)-\Pi_{k_C}m(P,\cdot)
\right\rangle_{L^2(\lambda)}.
\]
Cauchy--Schwarz and the approximation bound in the same cited statement yield
\[
\begin{aligned}
|\mathbb E_{P,n}\widehat C_n-C(P)|
&\le
\|e(P,\cdot)-\Pi_{k_C}e(P,\cdot)\|_2
\|m(P,\cdot)-\Pi_{k_C}m(P,\cdot)\|_2\\
&\le
25\left(\frac12\right)\left(\frac9{16}\right)
k_C^{-(\alpha+\beta)}
=\frac{225}{32}k_C^{-s}
\le8k_C^{-s}.
\end{aligned}
\]
Here \(0<\alpha\le1\), \(0<\beta<1\), and \(k_C\ge1\), so the approximation bound applies throughout the stated range.

Independence of the records and \(0\le AY\le1\) give
\[
\operatorname{Var}_{P,n}(\overline{AY}_n)
=\frac{\operatorname{Var}_P(AY)}n
\le\frac1n.
\]
Also, \cref{obj:lem:projection-control} supplies
\[
\operatorname{Var}_{P,n}\{\mathbb U_{n,k_C}(A,Y)\}
\le\frac4n+\frac{2k_C}{n(n-1)}.
\]
The variance identities for a sum and a difference give
\[
\begin{aligned}
\operatorname{Var}_{P,n}(\widehat C_n)
={}&2\operatorname{Var}_{P,n}(\overline{AY}_n)
+2\operatorname{Var}_{P,n}\{\mathbb U_{n,k_C}(A,Y)\}\\
&-\operatorname{Var}_{P,n}
 \{\overline{AY}_n+\mathbb U_{n,k_C}(A,Y)\}\\
\le{}&\frac2n+2\left\{\frac4n+\frac{2k_C}{n(n-1)}\right\}
=V_n(k_C),
\end{aligned}
\]
where the inequality uses nonnegativity of the last variance.
The finite-rank statistics are square integrable, so all these identities and subsequent expectations are well defined.

Cauchy--Schwarz on the probability space gives
\[
\mathbb E_{P,n}
\left|\widehat C_n-\mathbb E_{P,n}\widehat C_n\right|
\le
\sqrt{\operatorname{Var}_{P,n}(\widehat C_n)}.
\]
Consequently,
\[
\begin{aligned}
\mathbb E_{P,n}|\widehat C_n|
&\le
|\mathbb E_{P,n}\widehat C_n|
+\sqrt{\operatorname{Var}_{P,n}(\widehat C_n)}\\
&\le
|C(P)|+8k_C^{-s}+\sqrt{V_n(k_C)}
\le |C(P)|+b_C.
\end{aligned}
\]
For the last step, \(V_n(k_C)>0\) and
\(\sqrt{V_n(k_C)}\le\sqrt{20V_n(k_C)}\).

By \cref{obj:lem:observable-odds},
\[
C(P)=\{\exp(\theta(P))-1\}S(P),
\qquad S(P)\ge s_0.
\]
The product defining \(S(P)\) is at most one pointwise, so
\[
S(P)
=\int p_{10}(P,x)p_{01}(P,x)\,\lambda(dx)\le1,
\]
where the two conditional off-diagonal probabilities are
\[
p_{10}(P,x)=P(A=1,Y=0\mid X=x),
\qquad
p_{01}(P,x)=P(A=0,Y=1\mid X=x).
\]
Since \(|\theta(P)|\le1/2\), the exponential derivative is bounded by
\(\exp(1/2)\) on the segment between \(0\) and \(\theta(P)\). Therefore
\[
|\exp(\theta(P))-1|
\le\exp(1/2)|\theta(P)|,
\qquad
|C(P)|\le\exp(1/2)|\theta(P)|.
\]
In particular, for \(P\in\mathcal M_r(\alpha,\beta)\),
\[
\mathbb E_{P,n}|\widehat C_n|
\le\exp(1/2)r+b_C.
\]
% lean: CausalSmith.Stat.LogoddsLowsmoothFrontier.cHat_integral_abs_le
% lean: CausalSmith.Stat.LogoddsLowsmoothFrontier.covarianceCoordinate_abs_le_effect

\item We next bound the interval length on every sample. Use the confidence endpoints and clipped inversion of \cref{obj:def:upper-handle}:
\[
\begin{gathered}
c_\pm=\widehat C_n\pm b_C,
\qquad
s_\pm=\operatorname{clip}_{[s_0,1]}(\widehat S_n\pm b_S),\\
F(c_{\mathrm{rect}},s_{\mathrm{rect}})
=
\log\!\left(
\operatorname{clip}_{[\exp(-1/2),\exp(1/2)]}
\left(1+\frac{c_{\mathrm{rect}}}{s_{\mathrm{rect}}}\right)
\right)
\quad
(c_{\mathrm{rect}}\in\mathbb R,\ s_{\mathrm{rect}}>0),\\
L=\min\{F(c_-,s_-),F(c_-,s_+)\},
\qquad
R=\max\{F(c_+,s_-),F(c_+,s_+)\}.
\end{gathered}
\]
Clipping is a composition of minimum and maximum with fixed constants, each of which contracts absolute differences. Thus
\[
s_0\le s_-,s_+,
\qquad
|s_--s_+|
\le|(\widehat S_n-b_S)-(\widehat S_n+b_S)|
=2b_S.
\]

For real numbers satisfying
\(\exp(-1/2)\le v_{\log}\le w_{\log}\), the inequality
\(\log t\le t-1\) for \(t>0\) gives
\[
\log w_{\log}-\log v_{\log}
=
\log\frac{w_{\log}}{v_{\log}}
\le
\frac{w_{\log}-v_{\log}}{v_{\log}}
\le
\exp(1/2)(w_{\log}-v_{\log}).
\]
Interchanging the two arguments handles the reverse order. Applying this bound to the clipped arguments, and using contraction of the exponential-range clipping, gives
\[
\left|
F(c_{\mathrm{rect}},s_{\mathrm{rect}})
-
F(c_{\mathrm{rect}}',s_{\mathrm{rect}}')
\right|
\le
\exp(1/2)
\left|
\frac{c_{\mathrm{rect}}}{s_{\mathrm{rect}}}
-
\frac{c_{\mathrm{rect}}'}{s_{\mathrm{rect}}'}
\right|
\]
for all real numerators and positive denominators.

For the rectangle points with
\[
c_{\mathrm{rect}},c_{\mathrm{rect}}'\in[c_-,c_+],
\qquad
s_{\mathrm{rect}},s_{\mathrm{rect}}'\in\{s_-,s_+\},
\]
we have
\[
|c_{\mathrm{rect}}-c_{\mathrm{rect}}'|\le2b_C,
\qquad
|c_{\mathrm{rect}}'|\le|\widehat C_n|+b_C,
\qquad
|s_{\mathrm{rect}}-s_{\mathrm{rect}}'|\le2b_S.
\]
Inserting the cross quotient with numerator
\(c_{\mathrm{rect}}'\) and denominator \(s_{\mathrm{rect}}\) yields
\[
\begin{aligned}
\left|
\frac{c_{\mathrm{rect}}}{s_{\mathrm{rect}}}
-
\frac{c_{\mathrm{rect}}'}{s_{\mathrm{rect}}'}
\right|
&\le
\frac{|c_{\mathrm{rect}}-c_{\mathrm{rect}}'|}
     {s_{\mathrm{rect}}}
+
\frac{|c_{\mathrm{rect}}'|
      |s_{\mathrm{rect}}-s_{\mathrm{rect}}'|}
     {s_{\mathrm{rect}}s_{\mathrm{rect}}'}\\
&\le
\frac{2b_C}{s_0}
+
\frac{2b_S(|\widehat C_n|+b_C)}{s_0^2}.
\end{aligned}
\]
It follows that every pairing of an upper-numerator corner and a lower-numerator corner satisfies the same difference bound. Choosing the corners attaining the maximum \(R\) and minimum \(L\), including either choice in a tie, gives
\[
R-L
\le
\mathsf L_{\mathrm{sens}}b_C
+
\mathsf H_{\mathrm{sens}}b_S(|\widehat C_n|+b_C).
\]
By \cref{obj:lem:fixed-budget-realization}, the literal rational output has additional length at most \(2/n\). Hence, pointwise,
\[
\left|
I^{\mathrm{up}}_{n,\mathsf E,\mathsf N}(\alpha,\beta)(\mathcal O,U)
\right|
\le
\mathsf L_{\mathrm{sens}}b_C
+
\mathsf H_{\mathrm{sens}}b_S(|\widehat C_n|+b_C)
+\frac2n.
\]
% lean: CausalSmith.Stat.LogoddsLowsmoothFrontier.inversion_four_corner_width_le
% lean: CausalSmith.Stat.LogoddsLowsmoothFrontier.upperOutput_length_le

\item For \(P\in\mathcal M_r(\alpha,\beta)\), integrate the pointwise bound. The ranks and radii are fixed before sampling, and the numerator estimate is integrable. Therefore
\[
\begin{aligned}
&\mathbb E_{P,n}
\left|
I^{\mathrm{up}}_{n,\mathsf E,\mathsf N}(\alpha,\beta)(\mathcal O,U)
\right|\\
&\quad\le
\mathsf L_{\mathrm{sens}}b_C
+
\mathsf H_{\mathrm{sens}}b_S
 \{\mathbb E_{P,n}|\widehat C_n|+b_C\}
+\frac2n\\
&\quad\le
\mathsf L_{\mathrm{sens}}b_C
+
\mathsf H_{\mathrm{sens}}b_S
 \{\exp(1/2)r+2b_C\}
+\frac2n\\
&\quad\le K_0d_n(\alpha,\beta;r).
\end{aligned}
\]
The last inequality is the deterministic calibration established above.
% lean: CausalSmith.Stat.LogoddsLowsmoothFrontier.upperInterval_expectedLength_le_radii

\item To establish honesty, take any \(P\in\mathcal M(\alpha,\beta)\).
The numerator bias and variance bounds have already been established. For the denominator, the conditional probabilities defined above satisfy
\[
p_{10}(P,x)=e(P,x)\{1-\mu_1(P,x)\},
\qquad
p_{01}(P,x)=\{1-e(P,x)\}\mu_0(P,x).
\]
All their factors lie in \([0,1]\). Adding and subtracting the product with the first factor evaluated at \(x\) and the second at \(z\) gives
\[
\begin{aligned}
|p_{10}(P,x)-p_{10}(P,z)|
&\le |e(P,x)-e(P,z)|
   +|\mu_1(P,x)-\mu_1(P,z)|\\
&\le |x-z|^\beta,\\
|p_{01}(P,x)-p_{01}(P,z)|
&\le |e(P,x)-e(P,z)|
   +|\mu_0(P,x)-\mu_0(P,z)|\\
&\le |x-z|^\beta.
\end{aligned}
\]
Thus both continuous functions have Hölder seminorm at exponent \(\beta\) bounded by one.

The denominator marks \(A(1-Y)\) and \((1-A)Y\) have conditional means \(p_{10}(P,\cdot)\) and \(p_{01}(P,\cdot)\).
The expectation and residual identities in \cref{obj:lem:projection-control} therefore give
\[
\mathbb E_{P,n}\widehat S_n-S(P)
=
-\left\langle
p_{10}(P,\cdot)-\Pi_{k_S}p_{10}(P,\cdot),
p_{01}(P,\cdot)-\Pi_{k_S}p_{01}(P,\cdot)
\right\rangle_{L^2(\lambda)}.
\]
Cauchy--Schwarz and the cited approximation bound imply
\[
\begin{aligned}
|\mathbb E_{P,n}\widehat S_n-S(P)|
&\le
\|p_{10}(P,\cdot)-\Pi_{k_S}p_{10}(P,\cdot)\|_2
\|p_{01}(P,\cdot)-\Pi_{k_S}p_{01}(P,\cdot)\|_2\\
&\le25k_S^{-2\beta}.
\end{aligned}
\]
The variance bound in the same cited statement applies to these bounded marks:
\[
\operatorname{Var}_{P,n}(\widehat S_n)
\le\frac4n+\frac{2k_S}{n(n-1)}
\le V_n(k_S).
\]
Together, the four coordinate moment bounds are
\[
\begin{gathered}
|\mathbb E_{P,n}\widehat C_n-C(P)|\le8k_C^{-s},
\qquad
|\mathbb E_{P,n}\widehat S_n-S(P)|\le25k_S^{-2\beta},\\
\operatorname{Var}_{P,n}(\widehat C_n)\le V_n(k_C),
\qquad
\operatorname{Var}_{P,n}(\widehat S_n)\le V_n(k_S).
\end{gathered}
\]
% lean: CausalSmith.Stat.LogoddsLowsmoothFrontier.coordinateEstimates_bias_identities
% lean: CausalSmith.Stat.LogoddsLowsmoothFrontier.numerator_projection_bias_bound
% lean: CausalSmith.Stat.LogoddsLowsmoothFrontier.denominator_projection_bias_bound
% lean: CausalSmith.Stat.LogoddsLowsmoothFrontier.coordinateEstimates_variance_bounds

\item Define the coordinate failure events by
\[
\begin{aligned}
\mathcal B_C
&=\{(\mathcal O,U):|\widehat C_n-C(P)|>b_C\},\\
\mathcal B_S
&=\{(\mathcal O,U):|\widehat S_n-S(P)|>b_S\}.
\end{aligned}
\]
The triangle inequality and the bias bounds imply
\[
\begin{aligned}
\mathcal B_C
&\subseteq
\left\{
|\widehat C_n-\mathbb E_{P,n}\widehat C_n|
\ge\sqrt{20V_n(k_C)}
\right\},\\
\mathcal B_S
&\subseteq
\left\{
|\widehat S_n-\mathbb E_{P,n}\widehat S_n|
\ge\sqrt{20V_n(k_S)}
\right\}.
\end{aligned}
\]
Both variance radii are strictly positive. Integrating each squared centered coordinate over its displayed event gives
\[
\begin{aligned}
20V_n(k_C)\,
\mathbb P_{P,n}
\left\{
|\widehat C_n-\mathbb E_{P,n}\widehat C_n|
\ge\sqrt{20V_n(k_C)}
\right\}
&\le\operatorname{Var}_{P,n}(\widehat C_n),\\
20V_n(k_S)\,
\mathbb P_{P,n}
\left\{
|\widehat S_n-\mathbb E_{P,n}\widehat S_n|
\ge\sqrt{20V_n(k_S)}
\right\}
&\le\operatorname{Var}_{P,n}(\widehat S_n).
\end{aligned}
\]
Consequently,
\[
\mathbb P_{P,n}(\mathcal B_C)\le\frac1{20},
\qquad
\mathbb P_{P,n}(\mathcal B_S)\le\frac1{20},
\qquad
\mathbb P_{P,n}(\mathcal B_C\cup\mathcal B_S)\le\frac1{10}.
\]

Outside their union,
\[
C(P)\in[c_-,c_+],
\qquad
\widehat S_n-b_S\le S(P)\le\widehat S_n+b_S.
\]
We have already shown \(s_0\le S(P)\le1\).
Clipping is nondecreasing and fixes every point of \([s_0,1]\), so
\[
s_-\le S(P)\le s_+.
\]
For a fixed positive denominator, \(F\) is nondecreasing in its numerator.
For a fixed real numerator \(c_{\mathrm{rect}}\), its behavior as the positive denominator varies follows from
\[
\begin{cases}
\displaystyle
\frac{c_{\mathrm{rect}}}{s_+}
\le\frac{c_{\mathrm{rect}}}{s_{\mathrm{rect}}}
\le\frac{c_{\mathrm{rect}}}{s_-},
& c_{\mathrm{rect}}\ge0,\\[6pt]
\displaystyle
\frac{c_{\mathrm{rect}}}{s_-}
\le\frac{c_{\mathrm{rect}}}{s_{\mathrm{rect}}}
\le\frac{c_{\mathrm{rect}}}{s_+},
& c_{\mathrm{rect}}<0,
\end{cases}
\qquad
s_{\mathrm{rect}}\in[s_-,s_+].
\]
The clipping and logarithm preserve these inequalities. Applying them first to \(c_-\), then to \(c_+\), and using numerator monotonicity yields
\[
\begin{aligned}
L
&\le F(c_-,S(P))
\le F(C(P),S(P))\\
&\le F(c_+,S(P))
\le R.
\end{aligned}
\]
By \cref{obj:lem:observable-odds} and the effect envelope,
\[
1+\frac{C(P)}{S(P)}
=\exp(\theta(P))
\in[\exp(-1/2),\exp(1/2)],
\qquad
F(C(P),S(P))=\theta(P).
\]
Thus \(\theta(P)\in[L,R]\) outside \(\mathcal B_C\cup\mathcal B_S\).
The containment supplied by \cref{obj:lem:fixed-budget-realization} gives
\[
[L,R]\subseteq
I^{\mathrm{up}}_{n,\mathsf E,\mathsf N}(\alpha,\beta)(\mathcal O,U)
\]
on every sample. Hence
\[
\begin{aligned}
1
&\le
\mathbb P_{P,n}
\left\{
\theta(P)\in
I^{\mathrm{up}}_{n,\mathsf E,\mathsf N}(\alpha,\beta)(\mathcal O,U)
\right\}
+\mathbb P_{P,n}(\mathcal B_C\cup\mathcal B_S)\\
&\le
\mathbb P_{P,n}
\left\{
\theta(P)\in
I^{\mathrm{up}}_{n,\mathsf E,\mathsf N}(\alpha,\beta)(\mathcal O,U)
\right\}
+\frac1{10}.
\end{aligned}
\]
The coverage probability is therefore at least \(9/10\).
The Borel output property from the same cited statement gives measurable coverage events and interval length. Since \(P\) was arbitrary in the ambient model, this proves
\[
I^{\mathrm{up}}_{n,\mathsf E,\mathsf N}(\alpha,\beta)
\in\mathcal A_n(\alpha,\beta).
\]
% lean: CausalSmith.Stat.LogoddsLowsmoothFrontier.coordinate_error_probability_le
% lean: CausalSmith.Stat.LogoddsLowsmoothFrontier.upperInterval_coverage_of_coordinate_moments

\item Now let \(n\ge1\). If \(n=1\), the first branch of
\cref{obj:def:upper-handle} returns
\[
I^{\mathrm{up}}_{1,\mathsf E,\mathsf N}(\alpha,\beta)(\mathcal O,U)
=[-1/2,1/2]
\]
at every input. The effect envelope gives coverage one, and its length is one. For every \(r\in[0,1/2]\),
\[
d_1(\alpha,\beta;r)=1+r,
\qquad
\mathbb E_{P,1}
\left|
I^{\mathrm{up}}_{1,\mathsf E,\mathsf N}(\alpha,\beta)(\mathcal O,U)
\right|
=1
\le K_0(1+r).
\]
If \(n\ne1\), the integer assumption \(n\ge1\) implies \(n\ge2\), and the honesty and expected-length bounds established above apply. Thus the same positive constant \(K_0\) gives both assertions for every admissible engine, policy, exponent pair, sample size, radius, and law in the theorem.
% lean: CausalSmith.Stat.LogoddsLowsmoothFrontier.upperInterval_one_honest
% lean: CausalSmith.Stat.LogoddsLowsmoothFrontier.upperInterval_one_expectedLength_bound

\item Finally, every branch of the literal output construction either takes a nonempty intersection of a rational endpoint box with the full effect box or returns the full effect box. Each resulting box has ordered rational endpoints and is contained in \([-1/2,1/2]\); the output includes both endpoints. Consequently, for every sample and seed there exist \(l,u\in\mathbb Q\) such that
\[
I^{\mathrm{up}}_{n,\mathsf E,\mathsf N}(\alpha,\beta)(\mathcal O,U)
=[l,u],
\qquad
-\frac12\le l\le u\le\frac12.
\]
For \(n=1\), the full-box branch gives the asserted equality with
\([-1/2,1/2]\) at every sample and seed.
% lean: CausalSmith.Stat.LogoddsLowsmoothFrontier.upperOutput_rational_shape
% lean: CausalSmith.Stat.LogoddsLowsmoothFrontier.upperInterval_one_set
\end{enumerate}
\end{proof}

\begin{proof}[Proof of \cref{obj:thm:sharp-honest-frontier}]
\begin{enumerate}
\item Fix the joint calibration constants \(\varepsilon,M\) supplied by
\cref{obj:lem:calibrated-support}, and use the corresponding positive constants
\(\kappa,C_{\mathrm{mix}}\) from \cref{obj:lem:original-record-mixtures}. Define
\[
L_{\mathrm{mix}}
=
\left\lceil
\max\{1,\kappa^{-1},100C_{\mathrm{mix}}\varepsilon^4\}
\right\rceil,
\qquad
d_{\mathrm{ceil}}=(L_{\mathrm{mix}}+1)^{-1/2}.
\]
These are finite absolute constants, and the ceiling gives
\[
L_{\mathrm{mix}}\ge1,\qquad
\frac1{L_{\mathrm{mix}}}\le\kappa,\qquad
\frac{C_{\mathrm{mix}}\varepsilon^4}{L_{\mathrm{mix}}}
\le\frac1{100},\qquad
d_{\mathrm{ceil}}>0.
\]
For the argument below, write the worst expected length of an admissible procedure as
\[
\mathcal W_n(\alpha,\beta;r,I)
=
\sup_{P\in\mathcal M_r(\alpha,\beta)}
\mathbb E_{P^{\otimes n}\otimes\lambda}
|I(\mathcal O,U)|.
\]
We work throughout with \((\alpha,\beta)\in\mathcal D\), \(n\ge1\),
\(r\in[0,1/2]\), and \(I\in\mathcal A_n(\alpha,\beta)\), unless a narrower
range is specified.
% lean: CausalSmith.Stat.LogoddsLowsmoothFrontier.frontier_mixture_lower_rates

\item We first establish the rank estimates used in both mixture constructions.
For every real \(\gamma_{\mathrm{rank}}\in(0,1/2]\), define
\[
k_{\mathrm{low}}(\gamma_{\mathrm{rank}},n)
=
\left\lceil
L_{\mathrm{mix}}n^{2/(2\gamma_{\mathrm{rank}}+1)}
\right\rceil.
\]
Since \(n^{2/(2\gamma_{\mathrm{rank}}+1)}\ge1\), the ceiling inequalities yield
\[
1\le k_{\mathrm{low}}(\gamma_{\mathrm{rank}},n),\qquad
L_{\mathrm{mix}}n^{2/(2\gamma_{\mathrm{rank}}+1)}
\le k_{\mathrm{low}}(\gamma_{\mathrm{rank}},n)
\le
(L_{\mathrm{mix}}+1)n^{2/(2\gamma_{\mathrm{rank}}+1)}.
\]
The exponent \(2/(2\gamma_{\mathrm{rank}}+1)\) is at least one, so
\[
L_{\mathrm{mix}}n
\le
L_{\mathrm{mix}}n^{2/(2\gamma_{\mathrm{rank}}+1)}
\le k_{\mathrm{low}}(\gamma_{\mathrm{rank}},n),
\qquad
\frac{n}{k_{\mathrm{low}}(\gamma_{\mathrm{rank}},n)}
\le\frac1{L_{\mathrm{mix}}}.
\]
Moreover, \(L_{\mathrm{mix}}^{2\gamma_{\mathrm{rank}}+1}\ge L_{\mathrm{mix}}\).
Raising the lower rank bound to the positive power \(2\gamma_{\mathrm{rank}}+1\) therefore gives
\[
k_{\mathrm{low}}(\gamma_{\mathrm{rank}},n)^{2\gamma_{\mathrm{rank}}+1}
\ge
L_{\mathrm{mix}}^{2\gamma_{\mathrm{rank}}+1}n^2
\ge L_{\mathrm{mix}}n^2,
\]
and hence
\[
n^2k_{\mathrm{low}}(\gamma_{\mathrm{rank}},n)^{-(2\gamma_{\mathrm{rank}}+1)}
\le\frac1{L_{\mathrm{mix}}}.
\]
Finally, raising the upper rank bound to the negative power
\(-\gamma_{\mathrm{rank}}\), and using
\(\gamma_{\mathrm{rank}}\le1/2\), gives
\[
\begin{aligned}
k_{\mathrm{low}}(\gamma_{\mathrm{rank}},n)^{-\gamma_{\mathrm{rank}}}
&\ge
(L_{\mathrm{mix}}+1)^{-\gamma_{\mathrm{rank}}}
n^{-2\gamma_{\mathrm{rank}}/(2\gamma_{\mathrm{rank}}+1)}\\
&\ge
d_{\mathrm{ceil}}\,
n^{-2\gamma_{\mathrm{rank}}/(2\gamma_{\mathrm{rank}}+1)}.
\end{aligned}
\]
All these estimates include \(n=1\).
% lean: CausalSmith.Stat.LogoddsLowsmoothFrontier.lower_rank_ceiling_bounds
% lean: CausalSmith.Stat.LogoddsLowsmoothFrontier.lower_rank_budget
% lean: CausalSmith.Stat.LogoddsLowsmoothFrontier.lower_rank_separation

\item Suppose first that \(\alpha+\beta\le1/2\). Define the mixed rank and amplitudes by
\[
k_{\mathrm{mixed}}=k_{\mathrm{low}}(\alpha+\beta,n),
\qquad
\eta=\varepsilon k_{\mathrm{mixed}}^{-\alpha},
\qquad
\zeta=\varepsilon k_{\mathrm{mixed}}^{-\beta}.
\]
The domain conditions imply \(0<\alpha+\beta\le1/2\) and
\(\alpha,\beta>0\). Since \(k_{\mathrm{mixed}}\ge1\),
\[
0<\eta\le\varepsilon,\qquad 0<\zeta\le\varepsilon,
\qquad
\eta\zeta=\varepsilon^2k_{\mathrm{mixed}}^{-(\alpha+\beta)}.
\]
The rank estimates and the calibration of \(L_{\mathrm{mix}}\) give
\[
\frac{n}{k_{\mathrm{mixed}}}\le\kappa
\]
and, by the laws of real powers,
\[
\begin{aligned}
C_{\mathrm{mix}}\frac{n^2}{k_{\mathrm{mixed}}}\eta^2\zeta^2
&=
C_{\mathrm{mix}}\varepsilon^4
n^2k_{\mathrm{mixed}}^{-(2\alpha+2\beta+1)}\\
&\le
\frac{C_{\mathrm{mix}}\varepsilon^4}{L_{\mathrm{mix}}}
\le\frac1{100}.
\end{aligned}
\]
Thus \cref{obj:lem:original-record-mixtures}, applied to the mixed priors of
\cref{obj:def:continuous-calibrated-priors}, supplies
\[
\mathcal H^2(Q_0,Q_1)\le\frac1{100}.
\]

We give the testing-to-length argument needed here and for the fair priors.
Let \(Q_{\mathrm{eval}}\) and \(Q_{\mathrm{alt}}\) be uniform finite mixtures
of original-record sample laws. More precisely, for
\(\iota_{\mathrm{prior}}\in\{\mathrm{eval},\mathrm{alt}\}\), take an integer \(k_{\iota_{\mathrm{prior}}}\ge0\), define
\[
\Sigma_{\iota_{\mathrm{prior}}}=\{-1,1\}^{k_{\iota_{\mathrm{prior}}}+1},
\qquad
Q_{\iota_{\mathrm{prior}}}=\frac1{|\Sigma_{\iota_{\mathrm{prior}}}|}
\sum_{\sigma\in\Sigma_{\iota_{\mathrm{prior}}}}P_{\iota_{\mathrm{prior}},\sigma}^{\otimes n},
\]
and suppose that their constituents have common effects
\[
\theta(P_{\mathrm{eval},\sigma})=t_{\mathrm{eval}},
\qquad
\theta(P_{\mathrm{alt},\sigma})=t_{\mathrm{alt}}.
\]
Assume that every evaluation constituent belongs to
\(\mathcal M_r(\alpha,\beta)\), every alternative constituent belongs to
\(\mathcal M(\alpha,\beta)\), and all constituents have uniform covariate law.

Let \(\tau_{\mathrm{lab}}\) be the uniform probability law on the four binary
label pairs, and define the original-record reference and mixture densities by
\[
\tau_{\mathrm{lab}}
=\frac14\sum_{(a_{\mathrm{lab}},y_{\mathrm{lab}})\in\{0,1\}^2}
\delta_{(a_{\mathrm{lab}},y_{\mathrm{lab}})},
\qquad
R_{\mathrm{sample}}=(\lambda\otimes\tau_{\mathrm{lab}})^{\otimes n},
\qquad
q_{\iota_{\mathrm{prior}}}=\frac{dQ_{\iota_{\mathrm{prior}}}}{dR_{\mathrm{sample}}}.
\]
These densities are nonnegative and integrate to one. With total variation defined by
\[
\operatorname{TV}(Q_{\mathrm{eval}},Q_{\mathrm{alt}})
=
\frac12\int
|q_{\mathrm{eval}}-q_{\mathrm{alt}}|\,dR_{\mathrm{sample}},
\]
the squared Hellinger distance satisfies
\[
\mathcal H^2(Q_{\mathrm{eval}},Q_{\mathrm{alt}})
=
\int
(\sqrt{q_{\mathrm{eval}}}-\sqrt{q_{\mathrm{alt}}})^2\,dR_{\mathrm{sample}}
=
2-2\int\sqrt{q_{\mathrm{eval}}q_{\mathrm{alt}}}\,dR_{\mathrm{sample}}.
\]
Cauchy--Schwarz gives
\[
\int\sqrt{q_{\mathrm{eval}}q_{\mathrm{alt}}}\,dR_{\mathrm{sample}}\le1,
\qquad
\int(\sqrt{q_{\mathrm{eval}}}+\sqrt{q_{\mathrm{alt}}})^2
\,dR_{\mathrm{sample}}\le4.
\]
Factoring the density difference and applying Cauchy--Schwarz again yields
\[
\operatorname{TV}(Q_{\mathrm{eval}},Q_{\mathrm{alt}})
\le
\frac12
\sqrt{\mathcal H^2(Q_{\mathrm{eval}},Q_{\mathrm{alt}})}
\left[
\int(\sqrt{q_{\mathrm{eval}}}+\sqrt{q_{\mathrm{alt}}})^2
\,dR_{\mathrm{sample}}
\right]^{1/2}
\le
\sqrt{\mathcal H^2(Q_{\mathrm{eval}},Q_{\mathrm{alt}})}.
\]
The product densities with respect to
\(R_{\mathrm{sample}}\otimes\lambda\) are the same densities, independent of
the seed. Integrating over that seed proves
\[
\operatorname{TV}(Q_{\mathrm{eval}}\otimes\lambda,
Q_{\mathrm{alt}}\otimes\lambda)
=
\operatorname{TV}(Q_{\mathrm{eval}},Q_{\mathrm{alt}}).
\]
Consequently, a squared Hellinger bound of \(1/100\) gives total variation
at most \(1/10\) for the sample-and-seed experiments.

Define the measurable coverage events
\[
\mathcal E_{\mathrm{eval}}
=\{(\mathcal O,U):t_{\mathrm{eval}}\in I(\mathcal O,U)\},
\qquad
\mathcal E_{\mathrm{alt}}
=\{(\mathcal O,U):t_{\mathrm{alt}}\in I(\mathcal O,U)\}.
\]
Global honesty and finite averaging imply
\[
(Q_{\mathrm{eval}}\otimes\lambda)(\mathcal E_{\mathrm{eval}})\ge\frac9{10},
\qquad
(Q_{\mathrm{alt}}\otimes\lambda)(\mathcal E_{\mathrm{alt}})\ge\frac9{10}.
\]
Transferring the second event by total variation, then using that a union has
probability at most one, gives
\[
(Q_{\mathrm{eval}}\otimes\lambda)
(\mathcal E_{\mathrm{eval}}\cap\mathcal E_{\mathrm{alt}})
\ge\frac9{10}+\left(\frac9{10}-\frac1{10}\right)-1
=\frac7{10}.
\]
Connectedness of the returned interval gives the pointwise inequality
\[
|I(\mathcal O,U)|
\ge
|t_{\mathrm{alt}}-t_{\mathrm{eval}}|\,
\mathbf1_{\mathcal E_{\mathrm{eval}}\cap\mathcal E_{\mathrm{alt}}}
(\mathcal O,U).
\]
Integration and finite averaging therefore show
\[
\begin{aligned}
\frac7{10}|t_{\mathrm{alt}}-t_{\mathrm{eval}}|
&\le
\mathbb E_{Q_{\mathrm{eval}}\otimes\lambda}|I(\mathcal O,U)|\\
&=
\frac1{|\Sigma_{\mathrm{eval}}|}
\sum_{\sigma\in\Sigma_{\mathrm{eval}}}
\mathbb E_{P_{\mathrm{eval},\sigma}^{\otimes n}\otimes\lambda}
|I(\mathcal O,U)|\\
&\le\mathcal W_n(\alpha,\beta;r,I).
\end{aligned}
\]
The procedure's measurability and bounded interval length justify these
integrals. This argument includes independently randomized procedures.
% lean: CausalSmith.Stat.LogoddsLowsmoothFrontier.finitePrior_hellinger_length_lower

For the mixed priors, \cref{obj:lem:calibrated-support} supplies model
membership and the effects
\[
\theta(P_{0,\sigma})=0,\qquad
\theta(P_{1,\sigma})=32\eta\zeta.
\]
The null constituents belong to every evaluation slice because \(r\ge0\).
Applying the preceding reduction with \(Q_0\) as evaluation prior gives
\[
\mathcal W_n(\alpha,\beta;r,I)
\ge\frac7{10}\,32\eta\zeta
\ge
\frac7{10}\,32\varepsilon^2d_{\mathrm{ceil}}\,
n^{-2(\alpha+\beta)/(2\alpha+2\beta+1)}.
\]
Since the denominator is positive,
\[
\frac{2(\alpha+\beta)}{2\alpha+2\beta+1}\le\frac12
\quad\Longleftrightarrow\quad
\alpha+\beta\le\frac12.
\]
Thus, defining
\[
c_{\mathrm{mix}}=\frac7{10}\,32\varepsilon^2d_{\mathrm{ceil}}>0,
\]
we have proved
\[
c_{\mathrm{mix}}n^{-\mathsf a(\alpha,\beta)}
\le\mathcal W_n(\alpha,\beta;r,I)
\qquad(\alpha+\beta\le1/2).
\]
% lean: CausalSmith.Stat.LogoddsLowsmoothFrontier.mixed_prior_length_lower
% lean: CausalSmith.Stat.LogoddsLowsmoothFrontier.frontier_mixture_lower_rates

\item For the radius-dependent lower bound, suppose \(r>0\) and define
\[
k_{\mathrm{fair}}=k_{\mathrm{low}}(2\beta,n),
\qquad
\delta=\varepsilon k_{\mathrm{fair}}^{-\beta},
\qquad
t_{\mathrm{fair}}=\frac r2.
\]
Here \(0<2\beta<1/2\), \(0<\delta\le\varepsilon\le1/100\), and
\(t_{\mathrm{fair}}\in(0,1/4]\). The rank estimates imply
\[
\frac{n}{k_{\mathrm{fair}}}\le\kappa,
\qquad
C_{\mathrm{mix}}\frac{n^2}{k_{\mathrm{fair}}}\delta^4
=
C_{\mathrm{mix}}\varepsilon^4
n^2k_{\mathrm{fair}}^{-(4\beta+1)}
\le\frac1{100}.
\]
By \cref{obj:lem:original-record-mixtures},
\[
\mathcal H^2\!\left(
Q_{t_{\mathrm{fair}}},
P_{*,t_{\mathrm{fair}},\delta}^{\otimes n}
\right)\le\frac1{100}.
\]
The fair support assertion in \cref{obj:lem:calibrated-support} supplies
model membership and
\[
\theta(P_{t_{\mathrm{fair}},\sigma})=t_{\mathrm{fair}},
\qquad
\theta(P_{*,t_{\mathrm{fair}},\delta})
=\mathsf T(t_{\mathrm{fair}},\delta).
\]
The random constituents belong to the evaluation slice because
\(0<t_{\mathrm{fair}}\le r\). The comparator can be regarded as a constant
finite prior: averaging identical copies leaves its sample law unchanged.
By \cref{obj:lem:exact-calibrations},
\[
\frac{t_{\mathrm{fair}}}{2}
\le\mathsf T(t_{\mathrm{fair}},\delta)\le t_{\mathrm{fair}},
\qquad
t_{\mathrm{fair}}-\mathsf T(t_{\mathrm{fair}},\delta)
\ge3t_{\mathrm{fair}}\delta^2.
\]
Using the random fair prior as evaluation prior in the testing reduction,
and then the rank separation estimate at \(\gamma_{\mathrm{rank}}=2\beta\),
we obtain
\[
\begin{aligned}
\mathcal W_n(\alpha,\beta;r,I)
&\ge
\frac7{10}
\bigl[t_{\mathrm{fair}}-\mathsf T(t_{\mathrm{fair}},\delta)\bigr]\\
&\ge
\frac7{10}\,\frac32 r\varepsilon^2k_{\mathrm{fair}}^{-2\beta}\\
&\ge
\frac7{10}\,\frac32\varepsilon^2d_{\mathrm{ceil}}\,
r n^{-4\beta/(4\beta+1)}.
\end{aligned}
\]
Define
\[
c_{\mathrm{radius}}
=\frac7{10}\,\frac32\varepsilon^2d_{\mathrm{ceil}}>0.
\]
We have therefore proved
\[
c_{\mathrm{radius}}r n^{-\mathsf b(\beta)}
\le\mathcal W_n(\alpha,\beta;r,I)
\qquad(0<r\le1/2).
\]
% lean: CausalSmith.Stat.LogoddsLowsmoothFrontier.fair_prior_length_lower
% lean: CausalSmith.Stat.LogoddsLowsmoothFrontier.frontier_mixture_lower_rates

\item Choose the positive finite absolute upper constant supplied by
\cref{obj:thm:finite-honest-upper}, and set
\[
C=K_0.
\]
For every admissible engine and naming policy, that result supplies
\[
I^{\mathrm{up}}_{n,\mathsf E,\mathsf N}(\alpha,\beta)
\in\mathcal A_n(\alpha,\beta),
\]
together with
\[
\mathbb E_{P^{\otimes n}\otimes\lambda}
\left|I^{\mathrm{up}}_{n,\mathsf E,\mathsf N}(\alpha,\beta)(\mathcal O,U)\right|
\le C\rho_n(\alpha,\beta;r)
\qquad(P\in\mathcal M_r(\alpha,\beta)).
\]
Here the identification of the diagnostic envelope with \(\rho_n\) is
\cref{obj:def:frontier-handle}. Taking the supremum over the slice and using
the infimum definition in \cref{obj:def:length-objective} gives
\[
J_n(\alpha,\beta;r)
\le
\mathcal W_n\!\left(
\alpha,\beta;r,I^{\mathrm{up}}_{n,\mathsf E,\mathsf N}(\alpha,\beta)
\right)
\le C\rho_n(\alpha,\beta;r).
\]
The slice is nonempty: the law
\[
P_{\mathrm{null}}=\lambda\otimes\tau_{\mathrm{lab}}
\]
has
\[
e(P_{\mathrm{null}},x)=\mu_0(P_{\mathrm{null}},x)
=\mu_1(P_{\mathrm{null}},x)=\frac12,
\qquad
g(P_{\mathrm{null}},x)=\nu(P_{\mathrm{null}},x)
=\theta(P_{\mathrm{null}})=0,
\]
and hence satisfies the model conditions and belongs to every slice with
\(r\ge0\).

The upper result also supplies Borel measurability, global \(90\%\) coverage,
and the ordered rational closed output in the effect envelope. The
construction in \cref{obj:def:upper-handle} depends on the fixed engine,
public inputs, selected names, and original records, and therefore satisfies
\[
I^{\mathrm{up}}_{n,\mathsf E,\mathsf N}(\alpha,\beta)(\mathcal O,U)
=
I^{\mathrm{up}}_{n,\mathsf E,\mathsf N}(\alpha,\beta)(\mathcal O).
\]
Its input list makes it radius independent. For \(n\ge2\),
\cref{obj:lem:fixed-budget-realization} supplies successful termination
within the two prescribed budgets and finite rational evaluation. For
\(n=1\), the construction returns \([-1/2,1/2]\) directly.
% lean: CausalSmith.Stat.LogoddsLowsmoothFrontier.frontier_upper_comparisons

\item Define the common lower constants by
\[
c_{\mathrm{numerator}}
=\min\{c_{\mathrm{mix}},3/16\},
\qquad
c=\frac12\min\{c_{\mathrm{numerator}},c_{\mathrm{radius}}\}.
\]
Both constants are strictly positive. If \(\alpha+\beta\le1/2\), the mixed
bound gives
\[
c_{\mathrm{numerator}}n^{-\mathsf a(\alpha,\beta)}
\le
c_{\mathrm{mix}}n^{-\mathsf a(\alpha,\beta)}
\le\mathcal W_n(\alpha,\beta;r,I).
\]
If \(\alpha+\beta>1/2\), positivity of \(2\alpha+2\beta+1\) gives
\[
\frac12\le\frac{2(\alpha+\beta)}{2\alpha+2\beta+1},
\qquad
\mathsf a(\alpha,\beta)=\frac12.
\]
By \cref{obj:lem:parametric-null-floor} and the infimum definition of \(J_n\),
\[
c_{\mathrm{numerator}}n^{-\mathsf a(\alpha,\beta)}
\le
\frac3{16}n^{-1/2}
\le J_n(\alpha,\beta;r)
\le\mathcal W_n(\alpha,\beta;r,I).
\]
This proves the numerator bound throughout \(\mathcal D\), with equality
\(\alpha+\beta=1/2\) included in the first case.
% lean: CausalSmith.Stat.LogoddsLowsmoothFrontier.frontier_numerator_lower_of_rough

For the radius bound, \(r>0\) is covered above. At \(r=0\), nonnegativity of
interval length gives
\[
c_{\mathrm{radius}}\,0\,n^{-\mathsf b(\beta)}
=0\le\mathcal W_n(\alpha,\beta;0,I).
\]
Thus
\[
c_{\mathrm{radius}}r n^{-\mathsf b(\beta)}
\le\mathcal W_n(\alpha,\beta;r,I)
\qquad(0\le r\le1/2).
\]
% lean: CausalSmith.Stat.LogoddsLowsmoothFrontier.frontier_radius_lower_of_positive

Each summand of \(\rho_n\) is nonnegative. Reducing both coefficients to
their minimum and adding the two lower inequalities yields
\[
\begin{aligned}
c\rho_n(\alpha,\beta;r)
&=
\frac{\min\{c_{\mathrm{numerator}},c_{\mathrm{radius}}\}}2
\left[n^{-\mathsf a(\alpha,\beta)}
+r n^{-\mathsf b(\beta)}\right]\\
&\le
\frac12\left[
\mathcal W_n(\alpha,\beta;r,I)
+\mathcal W_n(\alpha,\beta;r,I)
\right]\\
&=\mathcal W_n(\alpha,\beta;r,I).
\end{aligned}
\]
This is the asserted converse for every honest procedure, including
independently randomized procedures.
% lean: CausalSmith.Stat.LogoddsLowsmoothFrontier.frontier_rate_lower_of_two

\item The honest class is nonempty: the constant procedure returning
\([-1/2,1/2]\) is Borel and covers every ambient-model effect by
\cref{obj:ass:effect-envelope}. Its length is one, while all procedure
objectives are nonnegative. Taking the infimum of the preceding
all-procedure lower bound is therefore legitimate and gives
\[
c\rho_n(\alpha,\beta;r)\le J_n(\alpha,\beta;r).
\]
Combining this with the upper comparison proves
\[
c\rho_n(\alpha,\beta;r)
\le J_n(\alpha,\beta;r)
\le
\mathcal W_n\!\left(
\alpha,\beta;r,I^{\mathrm{up}}_{n,\mathsf E,\mathsf N}(\alpha,\beta)
\right)
\le C\rho_n(\alpha,\beta;r).
\]
For the upper procedure, seed independence implies, for every law \(P\),
\[
\mathbb E_{P^{\otimes n}\otimes\lambda}
\left|I^{\mathrm{up}}_{n,\mathsf E,\mathsf N}(\alpha,\beta)(\mathcal O,U)\right|
=
\mathbb E_{P^{\otimes n}}
\left|I^{\mathrm{up}}_{n,\mathsf E,\mathsf N}(\alpha,\beta)(\mathcal O)\right|,
\]
which gives the displayed comparison in the statement. The constants were
chosen independently of the engine and naming policy.

By \cref{obj:lem:concrete-arithmetic-engine},
\(\mathsf E_0\) is admissible, and
\cref{obj:lem:fixed-budget-realization} supplies admissibility of
\(\mathsf N_0\). Specializing the universal comparison and operational
properties to this pair gives
\[
I_n^\star=I^{\mathrm{up}}_{n,\mathsf E_0,\mathsf N_0},
\qquad c_\star=c,\qquad C_\star=C,\qquad n_0=1.
\]
% lean: CausalSmith.Stat.LogoddsLowsmoothFrontier.universalFrontier_of_procedure_lower

\item Let \(\mathcal K\subseteq\mathcal D\) be nonempty and compact. The constant
functions with values \(c\) and \(1\) have singleton images on \(\mathcal K\), so
\[
\inf_{(\alpha,\beta)\in\mathcal K}c=c>0,
\qquad
\sup_{(\alpha,\beta)\in\mathcal K}1=1.
\]
The constant upper bound gives
\[
\sup_{(\alpha,\beta)\in\mathcal K}C\le C<\infty.
\]
These are the required compact-uniform conditions for the same constants
and threshold.
% lean: CausalSmith.Stat.LogoddsLowsmoothFrontier.frontier_compact_uniform

\item It remains to establish the exponent identities. The domain conditions give
\[
\alpha>0,\qquad
2\alpha+2\beta+1>0,\qquad
4\beta+1>0.
\]
As already verified by comparison of the two entries of the minimum,
\[
\mathsf a(\alpha,\beta)
=
\frac{2(\alpha+\beta)}{2\alpha+2\beta+1}
\quad\text{if }\alpha+\beta\le\frac12,
\qquad
\mathsf a(\alpha,\beta)=\frac12
\quad\text{if }\alpha+\beta\ge\frac12.
\]
In the first range, subtraction over the positive common denominator gives
\[
\begin{aligned}
\mathsf a(\alpha,\beta)-\mathsf b(\beta)
&=
\frac{
2(\alpha+\beta)(4\beta+1)
-4\beta(2\alpha+2\beta+1)}
{(2\alpha+2\beta+1)(4\beta+1)}\\
&=
\frac{2(\alpha-\beta)}
{(2\alpha+2\beta+1)(4\beta+1)}.
\end{aligned}
\]
In the second range,
\[
\mathsf a(\alpha,\beta)-\mathsf b(\beta)
=
\frac12-\frac{4\beta}{4\beta+1}
=
\frac{1-4\beta}{2(4\beta+1)}.
\]
If \(\alpha+\beta\le1/2\), positivity follows from \(\alpha>\beta\) and
the positive denominators. If \(\alpha+\beta>1/2\), it follows from
\(\beta<1/4\) and the positive denominator. Thus
\[
\mathsf a(\alpha,\beta)-\mathsf b(\beta)>0
\qquad((\alpha,\beta)\in\mathcal D).
\]
At \(\alpha+\beta=1/2\), both branch identities apply to the same difference,
and consequently
\[
\frac{2(\alpha-\beta)}
{(2\alpha+2\beta+1)(4\beta+1)}
=
\mathsf a(\alpha,\beta)-\mathsf b(\beta)
=
\frac{1-4\beta}{2(4\beta+1)}.
\]
This establishes agreement at the boundary and completes the proof.
% lean: CausalSmith.Stat.LogoddsLowsmoothFrontier.frontier_elbow_identities
\end{enumerate}
\end{proof}

\begin{proof}[Proof of \cref{obj:thm:full-honest-frontier-solution}]
\begin{enumerate}
\item Choose the absolute constants \(c,C\) supplied by
\cref{obj:thm:sharp-honest-frontier}, so that \(c>0\) and \(C>0\).
For every admissible engine \(\mathsf E\), admissible naming policy
\(\mathsf N\), exponent pair \((\alpha,\beta)\in\mathcal D\), and integer
\(n\ge1\), that result applies to the literal construction
\[
I^\star_{n,\mathsf E,\mathsf N}(\alpha,\beta)
=
I^{\mathrm{up}}_{n,\mathsf E,\mathsf N}(\alpha,\beta).
\]
It supplies membership in \(\mathcal A_n(\alpha,\beta)\), Borel
measurability, and termination on every valid named input within the
prescribed budgets. The inputs are the fixed engine, the public exponents
and their selected names, and the original observed records and their
selected covariate names, as specified in
\cref{obj:def:resolutions,obj:def:upper-handle}. For every sample and seed,
the returned interval has the form
\[
[l,u],\qquad l,u\in\mathbb Q,\qquad
-\frac12\le l\le u\le\frac12.
\]
For every \(r\in[0,1/2]\), the same result gives
\[
c\rho_n(\alpha,\beta;r)
\le J_n(\alpha,\beta;r)
\le
\sup_{P\in\mathcal M_r(\alpha,\beta)}
E_{P^{\otimes n}\otimes\lambda}
\left[
\left|
I^{\mathrm{up}}_{n,\mathsf E,\mathsf N}(\alpha,\beta)
(\mathcal O,U)
\right|
\right]
\le C\rho_n(\alpha,\beta;r),
\]
where \(\lambda\) is the uniform law of the independent seed \(U\).
The construction uses the same inputs for every evaluation radius.

Retain the choices
\[
c_\star(\alpha,\beta)=c,\qquad
C_\star(\alpha,\beta)=C,\qquad
n_0(\alpha,\beta)=1.
\]
The compact-uniform conclusion of
\cref{obj:thm:sharp-honest-frontier} states that, for every nonempty compact
\(\mathcal K\subset\mathcal D\),
\[
0<\inf_{(\alpha,\beta)\in\mathcal K}c_\star(\alpha,\beta),
\qquad
\sup_{(\alpha,\beta)\in\mathcal K}C_\star(\alpha,\beta)<\infty,
\qquad
\sup_{(\alpha,\beta)\in\mathcal K}n_0(\alpha,\beta)=1.
\]
Its exponent conclusions give, throughout \(\mathcal D\),
\[
\mathsf a(\alpha,\beta)-\mathsf b(\beta)>0,
\]
\[
\mathsf a(\alpha,\beta)-\mathsf b(\beta)=
\begin{cases}
\displaystyle
\frac{2(\alpha-\beta)}
{(2\alpha+2\beta+1)(4\beta+1)},
&\alpha+\beta\le1/2,\\[6pt]
\displaystyle
\frac{1-4\beta}{2(4\beta+1)},
&\alpha+\beta\ge1/2,
\end{cases}
\]
and, when \(\alpha+\beta=1/2\),
\[
\frac{2(\alpha-\beta)}
{(2\alpha+2\beta+1)(4\beta+1)}
=
\frac{1-4\beta}{2(4\beta+1)}.
\]
Finally, its procedure-wise converse supplies, for every
\((\alpha,\beta)\in\mathcal D\), \(r\in[0,1/2]\), \(n\ge1\), and
\(I\in\mathcal A_n(\alpha,\beta)\),
\[
c\rho_n(\alpha,\beta;r)
\le
\sup_{P\in\mathcal M_r(\alpha,\beta)}
E_{P^{\otimes n}\otimes\lambda}
\bigl[|I(\mathcal O,U)|\bigr],
\]
including procedures using independent randomization.
% lean: CausalSmith.Stat.LogoddsLowsmoothFrontier.sharp_honest_frontier

\item By \cref{obj:lem:concrete-arithmetic-engine}, the concrete engine
\(\mathsf E_0\) is admissible. By the first assertion of
\cref{obj:lem:fixed-budget-realization}, the dyadic-floor naming policy
\(\mathsf N_0\) is admissible. Specialize the universal conclusion above
to this pair. The resulting procedure is exactly
\[
I_n^\star(\alpha,\beta)
=
I^{\mathrm{up}}_{n,\mathsf E_0,\mathsf N_0}(\alpha,\beta),
\]
as defined in \cref{obj:def:frontier-handle}. Hence, for every
\((\alpha,\beta)\in\mathcal D\) and \(n\ge1\),
\[
I_n^\star(\alpha,\beta)\in\mathcal A_n(\alpha,\beta),
\]
and every sample and seed produces a closed interval with ordered
rational endpoints in \([-1/2,1/2]\). For every \(r\in[0,1/2]\),
\[
c\rho_n(\alpha,\beta;r)
\le J_n(\alpha,\beta;r)
\le
\sup_{P\in\mathcal M_r(\alpha,\beta)}
E_{P^{\otimes n}\otimes\lambda}
\bigl[|I_n^\star(\alpha,\beta)(\mathcal O,U)|\bigr]
\le C\rho_n(\alpha,\beta;r).
\]
This specialization preserves the constants chosen above and completes
the concrete attainment assertion together with the universal assertions.
% lean: CausalSmith.Stat.LogoddsLowsmoothFrontier.full_honest_frontier_solution
\end{enumerate}
\end{proof}
